\section{An explicit three-dimensional viscous strain layer}
\label{sec:viscous-strain}

The preceding phase-amplitude construction has a directly computable
three-dimensional counterpart. This section derives that counterpart in the
Cartesian Navier--Stokes equation, with its pressure and viscosity intact,
and places any finite interval of its central dynamics inside a smooth,
compactly supported forced solution. This is a local building block. The
constructed compact solution is global and therefore establishes no
nonexistence alternative from Section~\ref{sec:target}.

\subsection{The full velocity and its exact evolution}

Fix $\sigma>0$, $\nu>0$, $a_0\in\R$, and
$(K_1,K_3)\in\R^2\setminus\{(0,0)\}$. Coordinates remain
$x=(x_1,x_2,x_3)\in\R^3$ and time is $t\geq0$. Define
\begin{align}
 k_1(t)&=K_1e^{-\sigma t},& k_3(t)&=K_3,&
 k(t)&=(k_1(t),0,K_3),\\
 q(t)&=K_1^2e^{-2\sigma t}+K_3^2,&
 \varphi(x,t)&=k_1(t)x_1+K_3x_3.
\end{align}
\begin{equation}
 a(t)=a_0\exp\left\{\sigma t-
 \nu\left[\frac{K_1^2}{2\sigma}(1-e^{-2\sigma t})+K_3^2t\right]\right\}.
 \label{eq:vs-amplitude}
\end{equation}
In particular $q(t)>0$ at each finite time. Put
\begin{equation}\label{eq:vs-fields}
 u(x,t)=\bigl(\sigma x_1,-\sigma x_2+a(t)\cos\varphi(x,t),0\bigr),
 \qquad p_0(x,t)=-\frac{\sigma^2}{2}(x_1^2+x_2^2).
\end{equation}

\begin{proposition}[Exact viscous wave over a strain]
The fields \eqref{eq:vs-fields} solve
\[
 u_t+(u\cdot\nabla)u+\nabla p_0=\nu\Delta u,
 \qquad \operatorname{div}u=0
\]
at every $(x,t)\in\R^3\times[0,\infty)$. They are smooth on this
entire time interval. Their vorticity and its supremum norm are exactly
\begin{equation}\label{eq:vs-vorticity}
 \omega=a(t)\sin\varphi\,(K_3,0,-k_1(t)),
 \qquad \|\omega(t)\|_\infty=|a(t)|\sqrt{q(t)}.
\end{equation}
The full velocity has infinite spatial energy at every time.
\end{proposition}
\begin{proof}
Write $D=\operatorname{diag}(\sigma,-\sigma,0)$,
$u_b=Dx$, and $v=a\cos\varphi\,e_2$. Direct differentiation gives
\[
 k'=-D^Tk,\quad (\partial_t+u_b\cdot\nabla)\varphi=0,
 \quad a'=(\sigma-\nu q)a.
\]
The second component $v$ is independent of $x_2$, so
$(v\cdot\nabla)v=0$ and $\operatorname{div}v=0$.
Moreover $(v\cdot\nabla)u_b=Dv=-\sigma v$ and
$\Delta v=-qv$. Thus
\[
 v_t+(u_b\cdot\nabla)v+(v\cdot\nabla)u_b
 =(a'-\sigma a)\cos\varphi\,e_2=-\nu qv.
\]
The base equations hold because
$(u_b\cdot\nabla)u_b=D^2x=(\sigma^2x_1,\sigma^2x_2,0)$,
$\nabla p_0=-D^2x$, $\Delta u_b=0$, and $\tr D=0$.
Adding these identities proves every momentum component and divergence.
Taking the Cartesian curl gives \eqref{eq:vs-vorticity}. Since
$k(t)\ne0$, the linear map $x\mapsto k(t)\cdot x$ takes every real
value, so $|\sin\varphi|$ attains one; this proves the supremum formula.
Every displayed coefficient is smooth for finite $t$, proving global
smoothness. Finally $|u(x,t)|^2\geq\sigma^2x_1^2$; its integral over
$[1,2]\times\R\times[0,1]$ is infinite, proving the energy assertion.
\end{proof}

For $a_0\ne0$ the exact logarithmic vorticity rate is
\begin{equation}\label{eq:vs-growth}
 \frac{d}{dt}\log\|\omega(t)\|_\infty
 =\sigma-\nu q-\frac{\sigma k_1(t)^2}{q}
 =\frac{\sigma K_3^2}{q}-\nu q.
\end{equation}
Indeed $q'=-2\sigma k_1^2$ and differentiation of
$\log|a|+\tfrac12\log q$ gives the displayed equality. Consequently
vorticity increases at a particular time exactly when
$\sigma K_3^2>\nu q(t)^2$. In the retained subfamily $K_1=0$,
$K_3\ne0$, the full formula becomes
\[
 \|\omega(t)\|_\infty
 =|a_0K_3|e^{(\sigma-\nu K_3^2)t}.
\]
This grows when $\sigma>\nu K_3^2$, but has no finite-time
singularity. In the different subfamily $K_3=0$, $K_1\ne0$,
\[
 \|\omega(t)\|_\infty
 =|a_0K_1|\exp\left[-\frac{\nu K_1^2}{2\sigma}
 (1-e^{-2\sigma t})\right].
\]
Here the changing wavevector cancels the amplitude's strain growth. These
two exact expressions keep the distinction between amplitude, spatial
frequency, and vorticity visible in a common family.

\subsection{A compact realization with an unchanged interior}

Fix $0<R_0<R_1$ and $0<T_0<T_1$. The following explicit cutoffs specify
the entire constructed data and force. Define
\[
 b(s)=\begin{cases}e^{-1/s},&s>0,\\0,&s\leq0,\end{cases}
 \qquad H(s)=\frac{b(s)}{b(s)+b(1-s)},
\]
and
\begin{equation}\label{eq:vs-cutoffs}
 \chi(x)=H\left(\frac{R_1^2-|x|^2}{R_1^2-R_0^2}\right),
 \qquad
 \eta(t)=H\left(\frac{T_1-t}{T_1-T_0}\right).
\end{equation}
The denominator in $H$ is positive for all real $s$: if $s\leq0$,
then $1-s>0$; if $s\geq1$, then $s>0$; and both are positive when
$0<s<1$. Each positive-side derivative of $b$ is $e^{-1/s}$ times
a polynomial in $1/s$, as induction under differentiation shows.
For each integer $m\geq0$, $z^me^{-z}\to0$ as $z\to\infty$:
the series for $e^{z/2}$ bounds $z^m e^{-z/2}$ by a constant, leaving
an additional factor $e^{-z/2}\to0$. Thus all these derivatives
vanish at zero and $b,H,\chi,\eta$ are smooth. In particular
$0\leq\chi,\eta\leq1$, $\chi=1$ on $|x|\leq R_0$, $\chi=0$
on $|x|\geq R_1$, $\eta=1$ for $0\leq t\leq T_0$, and
$\eta=0$ for $t\geq T_1$.

An explicit vector potential for \eqref{eq:vs-fields} is
\begin{equation}\label{eq:vs-potential}
 A(x,t)=-\frac13x\times(Dx)
       +\frac{a(t)}{q(t)}(e_2\times k(t))\sin\varphi(x,t).
\end{equation}
Here the curl orientation is the original right-handed Cartesian one.
To check the base term, the product identity
\[
 \curl(X\times Y)=X\operatorname{div}Y
 -Y\operatorname{div}X+(Y\cdot\nabla)X-(X\cdot\nabla)Y
\]
follows by expanding its three components. For $X=x,Y=Dx$, its
right side is $0-3Dx+Dx-Dx=-3Dx$. For the wave term,
$k\times(e_2\times k)=q e_2$ since $k\cdot e_2=0$.
Taking its curl therefore gives $a\cos\varphi\,e_2$.
Together these calculations prove $\curl A=u$.

Set
\begin{equation}\label{eq:vs-compact-fields}
 w=\nabla\chi\times A,\qquad V=\chi u+w=\curl(\chi A),
 \qquad U=\eta V,\qquad P=\eta^2\chi p_0.
\end{equation}
Prescribe the force by the following expanded, finite expression, where
$\nabla\chi\cdot\nabla u=\sum_{j=1}^3(\partial_j\chi)\partial_j u$:
\begin{align}
 F={}&\eta'V+\eta\partial_t w
 -\nu\eta\bigl[2\nabla\chi\cdot\nabla u+(\Delta\chi)u+\Delta w\bigr]
 \nonumber\\
 &+\eta\chi(\eta\chi-1)(u\cdot\nabla)u
 +\eta\chi(\eta-1)\nabla p_0\nonumber\\
 &+\eta^2\bigl[\chi(u\cdot\nabla\chi)u
       +\chi(u\cdot\nabla)w+\chi(w\cdot\nabla)u
 \nonumber\\
 &\hspace{23mm}+(w\cdot\nabla\chi)u+(w\cdot\nabla)w
       +p_0\nabla\chi\bigr].
 \label{eq:vs-compact-force}
\end{align}
All factors in this formula have been explicitly defined, including their
initial values and cutoffs. In particular no solution-dependent unknown
force is left to choose.

\begin{proposition}[Exact localization and admissibility]
The fields $(U,P,F)$ in
\eqref{eq:vs-compact-fields}--\eqref{eq:vs-compact-force} satisfy
\[
 U_t+(U\cdot\nabla)U+\nabla P=\nu\Delta U+F,
 \qquad \operatorname{div}U=0.
\]
They are smooth for all $t\geq0$. Their spatial supports lie in
$\{|x|\leq R_1\}$, and all three vanish for $t\geq T_1$.
The initial datum $U(x,0)=\curl(\chi A(x,0))$ satisfies
\eqref{eq:schwartz-data}, the force satisfies
\eqref{eq:schwartz-force}, and $U$ satisfies \eqref{eq:energy}.
On $|x|<R_0$, $0\leq t\leq T_0$, one has exactly
$(U,P,F)=(u,p_0,0)$.
\end{proposition}
\begin{proof}
The divergence of a curl is zero because each mixed second derivative
appears twice with opposite signs. Since $\eta$ is spatially constant,
this proves $\operatorname{div}U=0$. Direct products give
\begin{align*}
 \Delta V&=\chi\Delta u+2\nabla\chi\cdot\nabla u
              +(\Delta\chi)u+\Delta w,\\
 (V\cdot\nabla)V
 &=\chi^2(u\cdot\nabla)u+\chi(u\cdot\nabla\chi)u
   +\chi(u\cdot\nabla)w+\chi(w\cdot\nabla)u\\
 &\qquad +(w\cdot\nabla\chi)u+(w\cdot\nabla)w,\\
 \nabla(\chi p_0)&=\chi\nabla p_0+p_0\nabla\chi.
\end{align*}
Insert these three identities into
$U_t+(U\cdot\nabla)U+\nabla P-\nu\Delta U$.
The terms $\eta\chi(u_t-\nu\Delta u)$ equal
$-\eta\chi[(u\cdot\nabla)u+\nabla p_0]$ by the already proved
uncut equation. Combining their two coefficients with the displayed
convection and pressure expansions produces
$\eta\chi(\eta\chi-1)$ and $\eta\chi(\eta-1)$ respectively.
Every remaining term is exactly \eqref{eq:vs-compact-force}, proving
the momentum equation with the prescribed force.

Smoothness follows from the explicit formulas, $q>0$ at finite time,
and flatness of the cutoffs at their endpoints. Outside the spatial
cutoff, $\chi$ and all its derivatives vanish, hence $w,V,U,P,F$
vanish. For $t\geq T_1$, $\eta$ and every time derivative vanish,
so $U,P,F$ vanish as well. Each mixed derivative of $F$ is consequently
bounded on the compact set $\{|x|\leq R_1,0\leq t\leq T_1\}$,
say by $M_{\alpha m}$. For every $K\in\mathbb N_0$, the choice
\[
 C_{\alpha mK}=M_{\alpha m}(1+R_1+T_1)^K
\]
proves \eqref{eq:schwartz-force}, including the exterior where the
derivative is zero. The same argument at $t=0$ proves the data condition.

For a quantitative energy bound define
\[
 a_* =|a_0|e^{\sigma T_1},\qquad
 k_* =\sqrt{K_1^2e^{-2\sigma T_1}+K_3^2}>0,
 \qquad c_\chi=\|\nabla\chi\|_\infty.
\]
The viscosity exponent in \eqref{eq:vs-amplitude} is nonpositive for
$0\leq t\leq T_1$, so $|a|\leq a_*$ there. Also
$\sqrt q\geq k_*$, $|Dx|\leq\sigma R_1$ on the spatial support,
and \eqref{eq:vs-potential} gives
$|A|\leq\sigma R_1^2/3+a_*/k_*$. Consequently
\[
 \int_{\R^3}|U(x,t)|^2\,dx
 \leq\frac{4\pi R_1^3}{3}
 \left[\sigma R_1+a_*+
 c_\chi\left(\frac{\sigma R_1^2}{3}+\frac{a_*}{k_*}\right)\right]^2
 \quad(t\geq0).
\]
Adding one to this finite right side provides the strict bound in
\eqref{eq:energy}. Finally, on the stated interior plateau
$\chi=\eta=1$, $w=0$ and all cutoff derivatives vanish; direct
substitution proves $(U,P,F)=(u,p_0,0)$ there.
\end{proof}

The original strain, both initial wavevector components, amplitude,
viscosity, inner and outer spatial radii, and both time endpoints have
all been retained. Localization supplies an actual map from this
nondecaying exact solution to admissible compact data and forcing with an
identical finite interior evolution. Increasing an amplification parameter
creates a different smooth solution; the displayed family alone does not
produce a finite-time singular solution for fixed prescribed data.
